% ==============================================================================
% T0/FFGFT-Theorie: Dokument 231
% Titel: Erweiterungen des Hilbertraum-Formalismus für volle FFGFT
% Autor: Johann Pascher
% Datum: Mai 2026
% Bezug: Dokumente 175 (Qubit-Zustandsräume), 174 (Spin-Qubit), 210 (Wicklungen),
%        202 (Feldtheorie), 230 (Hilbertraum-Übersetzung)
% ==============================================================================

\section*{Vorbemerkung}

Dokument~230 hat gezeigt, wie sich FFGFT-Aussagen \emph{in den
Standard-Hilbertraum übersetzen lassen} -- mit Reduktionen, die bei
jeder Stufe Information verlieren. Dieses Dokument dreht die
Frage um:

\begin{quote}
\itshape
\enquote{Welche zusätzlichen mathematischen Strukturen muss man dem
Standard-Hilbertraum hinzufügen, damit er die volle FFGFT-Geometrie
\textbf{nativ} -- also ohne Reduktion -- abbilden kann?}
\end{quote}

Die Antwort hat eine bemerkenswerte Eigenschaft: jede notwendige
Erweiterung entspricht einer \emph{etablierten mathematischen
Struktur} aus anderen Bereichen der theoretischen Physik. FFGFT
\enquote{erfindet} keine neue Mathematik; sie ist die spezifische
Kombination etablierter Erweiterungen, die parameterfrei alle
Standardmodell-Konstanten liefert.

% ==============================================================================
\section*{Zeichenerklärung}
% ==============================================================================

\subsection*{Geometrische Räume und Hilberträume}

\begin{center}
\renewcommand{\arraystretch}{1.3}
{\small
\adjustbox{max width=\textwidth}{%
\begin{tabular}{>{\raggedright\arraybackslash}p{2.8cm}>{\raggedright\arraybackslash}p{8.2cm}}
\toprule
\textbf{Symbol} & \textbf{Bedeutung} \\
\midrule
$\mathbb{R}$ & reelle Zahlen \\
$\mathbb{C}$ & komplexe Zahlen \\
$\mathbb{C}^2$ & Standard-Qubit-Hilbertraum \\
$\mathbb{Z}$ & ganze Zahlen (Wicklungsindex) \\
$\aleph_0$ & abzählbar unendlich (Hilbertraum-Dimension nach Erweiterung 2) \\
$S^1, S^2$ & Kreis, 2-Sphäre \\
$T^n$ & $n$-Torus, $n$ kompakte Schleifen \\
$L^2(M, S)$ & Schnitte des Spinorbündels $S$ über Mannigfaltigkeit $M$ \\
$\mathcal{H}_{\text{std}}$ & Standard-Hilbertraum \\
$\mathcal{H}^{T^n}_{\text{FFGFT}}$ & FFGFT-Hilbertraum auf $T^n$ \\
$\bigoplus_n$ & direkte Summe über Index $n$ \\
$\bigotimes$ & Tensorprodukt \\
\bottomrule
\end{tabular}%
}
}
\renewcommand{\arraystretch}{1.0}
\end{center}

\subsection*{Algebraische Strukturen}

\begin{center}
\renewcommand{\arraystretch}{1.3}
{\small
\adjustbox{max width=\textwidth}{%
\begin{tabular}{>{\raggedright\arraybackslash}p{2.8cm}>{\raggedright\arraybackslash}p{8.2cm}}
\toprule
\textbf{Symbol} & \textbf{Bedeutung} \\
\midrule
$SU(2)$ & spezielle unitäre Gruppe (Standard) \\
$SU_q(2)$ & $q$-deformierte Quantengruppe (Drinfeld-Jimbo) \\
$q$ & Deformationsparameter, $q = e^{i\pi(1-K_{\text{frak}})} = e^{i\pi/75}$ \\
$[n]_q$ & $q$-deformierte Zahl, $[n]_q = (q^n - q^{-n})/(q - q^{-1})$ \\
$\sigma_x, \sigma_y, \sigma_z$ & Pauli-Matrizen \\
$\sigma_x^{\text{FFGFT}}$ & FFGFT-modifiziertes $\sigma_x = K_{\text{frak}} \sigma_x$ \\
$[A, B]$ & Kommutator $AB - BA$ \\
$\{A, B\}$ & Antikommutator $AB + BA$ \\
$\mathbb{1}$ & Einheitsmatrix \\
$\delta_{ij}, \epsilon_{ijk}$ & Kronecker-Delta, Levi-Civita-Tensor \\
\bottomrule
\end{tabular}%
}
}
\renewcommand{\arraystretch}{1.0}
\end{center}

\subsection*{FFGFT-Geometrie-Parameter}

\begin{center}
\renewcommand{\arraystretch}{1.3}
{\small
\adjustbox{max width=\textwidth}{%
\begin{tabular}{>{\raggedright\arraybackslash}p{2.8cm}>{\raggedright\arraybackslash}p{8.2cm}}
\toprule
\textbf{Symbol} & \textbf{Bedeutung} \\
\midrule
$\xi_0 = 4/30000$ & dimensionsloser FFGFT-Geometrieparameter \\
$\xi$ & alternative Schreibweise für $\xi_0$ \\
$D_f = 2{,}999867$ & fraktale Raumzeitdimension \\
$K_{\text{frak}}$ & fraktale Korrektur, $K_{\text{frak}} = 1 - 100\xi_0 \approx 0{,}9867$ \\
$1 - K_{\text{frak}}$ & Deformationsphase, $\approx 0{,}0133$ \\
$v$ & Higgs-Vakuumerwartungswert, $\approx 246$ GeV \\
$\hbar, c$ & Plancksches Wirkungsquantum, Lichtgeschwindigkeit \\
\bottomrule
\end{tabular}%
}
}
\renewcommand{\arraystretch}{1.0}
\end{center}

\subsection*{Wicklungsstruktur (Erweiterungen 2 und 4)}

\begin{center}
\renewcommand{\arraystretch}{1.3}
{\small
\adjustbox{max width=\textwidth}{%
\begin{tabular}{>{\raggedright\arraybackslash}p{2.8cm}>{\raggedright\arraybackslash}p{8.2cm}}
\toprule
\textbf{Symbol} & \textbf{Bedeutung} \\
\midrule
$n \in \mathbb{Z}$ & zyklische Wicklungs-Quantenzahl (Schleife) \\
$|\psi; n\rangle$ & Zustand mit Wicklungszahl $n$ \\
$\mathbb{C}^2_n$ & Qubit-Hilbertraum im $n$-Wicklungssektor \\
$E_n$ & Energie der $n$-ten Wicklungsanregung \\
$R$ & Schleifenradius (Hauptradius des Torus) \\
$a^\dagger, a$ & Erzeugungs- und Vernichtungsoperatoren für Wicklungsanregungen \\
$k_t \in \mathbb{Z}$ & zeitliche Wicklungs-Quantenzahl \\
$f(k_t)$ & periodische Energiefunktion der $k_t$-Wicklung \\
$(n_x, n_y, n_z)$ & räumliche Wicklungs-Quantenzahlen \\
$(n_\theta, n_\phi)$ & Spin-Wicklungs-Quantenzahlen \\
\bottomrule
\end{tabular}%
}
}
\renewcommand{\arraystretch}{1.0}
\end{center}

\subsection*{Operatoren auf Faserbündeln (Erweiterung 3)}

\begin{center}
\renewcommand{\arraystretch}{1.3}
{\small
\adjustbox{max width=\textwidth}{%
\begin{tabular}{>{\raggedright\arraybackslash}p{2.8cm}>{\raggedright\arraybackslash}p{8.2cm}}
\toprule
\textbf{Symbol} & \textbf{Bedeutung} \\
\midrule
$\vec L$ & Bahn-Drehimpuls-Operator \\
$\vec S$ & Spin-Operator \\
$L_i, S_j$ & Komponenten von $\vec L, \vec S$ \\
$H_{\text{SO}}$ & Spin-Bahn-Kopplungs-Hamilton \\
$\xi(r)$ & radiale Spin-Bahn-Kopplungsfunktion \\
$Z$ & Kernladungszahl \\
$S$ & Spinorbündel über $T^3$ oder $T^4$ \\
$\mathcal{D}_{T^4}$ & Dirac-artiger Operator auf $T^4$ \\
\bottomrule
\end{tabular}%
}
}
\renewcommand{\arraystretch}{1.0}
\end{center}

\subsection*{Standard-QM-Größen}

\begin{center}
\renewcommand{\arraystretch}{1.3}
{\small
\adjustbox{max width=\textwidth}{%
\begin{tabular}{>{\raggedright\arraybackslash}p{2.8cm}>{\raggedright\arraybackslash}p{8.2cm}}
\toprule
\textbf{Symbol} & \textbf{Bedeutung} \\
\midrule
$|\psi\rangle$ & Quantenzustand \\
$|0\rangle, |1\rangle$ & Berechnungsbasis-Zustände \\
$\alpha, \beta$ & komplexe Amplituden \\
$\psi(x,y,z)$ & Ortswellenfunktion \\
$\omega = 2\pi/T$ & Kreisfrequenz \\
$E$ & Energie \\
$m, m_i$ & Masse, Masse des Fermions $i$ \\
$\alpha$ & Feinstrukturkonstante \\
$\nabla^2$ & Laplace-Operator \\
\bottomrule
\end{tabular}%
}
}
\renewcommand{\arraystretch}{1.0}
\end{center}

% ==============================================================================
\section*{Die vier notwendigen Erweiterungen}
% ==============================================================================

Aus Dok.~230 wissen wir: die Hilbertraum-Übersetzung eines einzelnen
Qubits passiert auf der Bloch-Sphäre ($S^2$), die durch zwei
Reduktionen aus der vollen FFGFT-Geometrie entsteht
($T^4 \to T^3 \to T^2 \to \text{Zylinder} \to S^2$). Jede dieser
Reduktionen entspricht dem Verlust einer Strukturschicht. Um die
volle FFGFT zu rekonstruieren, müssen wir vier Erweiterungen
zurückgewinnen.

\textbf{Übersicht:}

\begin{center}
\renewcommand{\arraystretch}{1.4}
{\small
\adjustbox{max width=\textwidth}{%
\begin{tabular}{>{\raggedright\arraybackslash}p{1.7cm}>{\raggedright\arraybackslash}p{4.0cm}>{\raggedright\arraybackslash}p{4.3cm}}
\toprule
\textbf{Stufe} & \textbf{Erweiterung} & \textbf{Mathematisches Vorbild} \\
\midrule
$S^2 \to$ Zylinder
  & deformierte SU(2)-Algebra
  & Quantengruppen $SU_q(2)$ \\[4pt]
Zyl.\ $\to T^2$
  & zusätzliche zyklische Quantenzahl
  & Kac-Moody-Algebren, Wicklungs-Zustände \\[4pt]
$T^2 \to T^3$
  & Bündelstruktur statt Tensorprodukt
  & Faserbündel, Kaluza-Klein \\[4pt]
$T^3 \to T^4$
  & zusätzliche zeitliche Wicklungsdimension
  & Kaluza-Klein, kompakte Extradimensionen \\
\bottomrule
\end{tabular}%
}
}
\renewcommand{\arraystretch}{1.0}
\end{center}

% ==============================================================================
\section{Erweiterung 1: deformierte SU(2)-Algebra}
% ==============================================================================

\subsection{Standard-Pauli-Algebra}

Die Pauli-Matrizen $\sigma_x, \sigma_y, \sigma_z$ erzeugen die
isotrope $SU(2)$-Algebra:
\begin{equation}
  [\sigma_i, \sigma_j] = 2i\epsilon_{ijk}\sigma_k,
  \qquad
  \{\sigma_i, \sigma_j\} = 2\delta_{ij}\mathbb{1}.
\end{equation}
\enquote{Isotrop} bedeutet hier: alle drei Achsen $x, y, z$ sind
gleichberechtigt; eine Drehung um $\sigma_x$ ist exakt äquivalent zu
einer Drehung um $\sigma_z$ nach geeigneter Basisänderung.

\subsection{FFGFT-Anisotropie}

In FFGFT gibt es einen \emph{ausgezeichneten Vorfaktor} für die
$\sigma_x$-Drehung:
\begin{equation}
  \sigma_x^{\text{FFGFT}} = K_{\text{frak}} \cdot \sigma_x,
  \qquad
  K_{\text{frak}} = 1 - 100\xi_0 \approx 0{,}9867.
\end{equation}
Die $\sigma_y$- und $\sigma_z$-Drehungen sind ungestört. Damit ist
die Algebra nicht mehr isotrop: $\sigma_x$ ist ausgezeichnet.

\subsection{Mathematische Identifikation: $SU_q(2)$}

Die $SU_q(2)$-Quantengruppe (Drinfeld-Jimbo, 1985) ist die
$q$-Deformation der $SU(2)$, mit modifizierten Vertauschungsrelationen.
Im Limes $q \to 1$ wird sie zur Standard-$SU(2)$. Die FFGFT-Anisotropie
entspricht der Identifikation:
\begin{equation}
  \boxed{
    q = e^{i\pi(1 - K_{\text{frak}})}
      = e^{i\pi \cdot 100\xi_0}
      = e^{i\pi/75}
      \approx 0{,}999 + 0{,}042i.
  }
\end{equation}

Die $q$-Phase ist klein ($\pi/75 \approx 0{,}042$ rad), aber
strukturell signifikant: sie bricht die kugelsymmetrische
$SU(2)$-Algebra zu einer anisotropen, achsenausgezeichneten Algebra.
Die Anisotropie ist die mathematische Erinnerung des Zylinders an
seine Torus-Herkunft.

\subsection{Was diese Erweiterung kostet}

\begin{itemize}[leftmargin=2em,itemsep=2pt,topsep=2pt]
  \item Hilbertraum bleibt $\mathbb{C}^2$ -- keine Dimensions-Erweiterung.
  \item Operatoren bleiben $2\times 2$-Matrizen -- keine
        Matrix-Erweiterung.
  \item Was sich ändert: die Strukturkonstanten der Algebra
        (Casimir-Operatoren, Multipliers).
  \item Konsequenz: alle $\pi$-Gatter weichen systematisch um
        $0{,}013\,\%$ vom idealen Wert ab (Dok.~175 §3, testbar in
        Hochpräzisions-Quantengatter-Experimenten).
\end{itemize}

% ==============================================================================
\section{\texorpdfstring{Erweiterung 2: zyklische Quantenzahl}{Erweiterung 2: zyklische Quantenzahl}}
% ==============================================================================

\subsection{Was beim Zylinder fehlt}

Der Zylinder $\mathbb{R} \times S^1$ hat \emph{eine} kompakte
Schleife (die Phase $\theta$) und \emph{eine} lineare Achse ($z$). Im
2-Torus $T^2 = S^1 \times S^1$ sind \emph{beide} Achsen kompakt --
das heißt, $z = +1$ und $z = -1$ sind durch eine zweite Schleife
verbunden. Im Bloch-Sphäre-Limes (alle Pole zusammenziehen) und
selbst im Zylinder-Limes (eine Schleife auseinanderziehen) verliert
man diese Verbindung.

\subsection{Hilbertraum-Erweiterung}

Um die zweite Schleife zu repräsentieren, muss der Hilbertraum eine
zusätzliche \emph{Wicklungs-Quantenzahl} $n \in \mathbb{Z}$ tragen.
Statt
\begin{equation}
  |\psi\rangle = \alpha|0\rangle + \beta|1\rangle \in \mathbb{C}^2
\end{equation}
wird der Zustand
\begin{equation}
  |\psi; n\rangle = \alpha|0; n\rangle + \beta|1; n\rangle,
\end{equation}
wobei $n$ angibt, wie oft die Mode-Konfiguration die zweite
Schleife des $T^2$ umläuft. Der Hilbertraum wird:
\begin{equation}
  \boxed{
    \mathcal{H}^{T^2}_{\text{FFGFT}}
    = \bigoplus_{n \in \mathbb{Z}} \mathbb{C}^2_n.
  }
\end{equation}
Das ist ein \emph{abzählbar unendlich-dimensionaler} Hilbertraum,
nicht mehr 2-dimensional.

\subsection{Niederenergielimit: warum wir das normalerweise nicht sehen}

Höhere Wicklungs-Anregungen ($n \neq 0$) haben höhere Energie:
\begin{equation}
  E_n \sim n^2 \cdot \frac{\hbar^2}{m R^2},
\end{equation}
wobei $R$ der Schleifenradius ist. Bei niedrigen Energien ist nur
$n = 0$ besetzt, und der Hilbertraum reduziert sich effektiv zu
$\mathbb{C}^2$ -- die Standard-Bloch-Sphäre. Höhere
$n$-Anregungen werden bei makroskopischen $R$ rasch unsichtbar.

\subsection{Mathematisches Vorbild}

Die Struktur $\bigoplus_n \mathbb{C}^2_n$ ist genau das, was in der
Stringtheorie als \emph{Wicklungs-Zustände kompakter Dimensionen}
auftritt, und in Kac-Moody-Algebren als \emph{Levels}. Auch hier ist
die Mathematik bekannt -- FFGFT importiert sie mit einer
spezifischen geometrischen Bedeutung.

\subsection{Was diese Erweiterung kostet}

\begin{itemize}[leftmargin=2em,itemsep=2pt,topsep=2pt]
  \item Hilbertraum-Dimension wird $\aleph_0$ statt $2$.
  \item Operatoren werden unendlich-dimensional, aber blockdiagonal:
        innerhalb jedes $n$-Sektors sind sie $2\times 2$.
  \item Zusätzliche Operatoren $a^\dagger, a$ propagieren zwischen
        Wicklungssektoren ($n \to n \pm 1$).
  \item In niederen Energien irrelevant; in der vollen Theorie
        notwendig (z.B.\ für Vakuum-Casimir-Effekte zwischen den
        Polen, Dok.~091/220).
\end{itemize}

% ==============================================================================
\section{Erweiterung 3: Bündelstruktur statt Tensorprodukt}
% ==============================================================================

\subsection{Standard-QM für ein Teilchen mit Spin}

Ein Spin-$\tfrac{1}{2}$-Teilchen im Raum hat in der Standard-QM den
Hilbertraum:
\begin{equation}
  \mathcal{H}_{\text{std}} = L^2(\mathbb{R}^3) \otimes \mathbb{C}^2.
\end{equation}
Das \emph{Tensorprodukt} bedeutet: Ortsraum und Spin-Raum sind
\textbf{unabhängige} Hilberträume. Operatoren des einen Faktors
kommutieren mit Operatoren des anderen:
\begin{equation}
  [L_i, S_j] = 0
  \qquad \text{(Bahn- und Spin-Drehimpuls vertauschen)}.
\end{equation}
Spin-Bahn-Kopplung muss als \emph{zusätzlicher Hamilton-Term}
eingeführt werden:
\begin{equation}
  H_{\text{SO}} = \xi(r)\, \vec L \cdot \vec S,
\end{equation}
wobei die \enquote{Spin-Bahn-Konstante} $\xi(r)$ ein freier
Parameter ist (typisch $\sim Z^4 \alpha^2$ für ein
wasserstoffartiges Atom).

\subsection{FFGFT: Spin und Ort auf demselben $T^3$}

In FFGFT entstehen sowohl der \enquote{Ort} als auch der
\enquote{Spin} aus Wicklungen auf demselben 3-Torus $T^3$ (Dok.~210):
\begin{itemize}[leftmargin=2em,itemsep=2pt,topsep=2pt]
  \item räumliche Wicklungen $(n_x, n_y, n_z)$: Modenraum-Position,
  \item Spin-Wicklungen $(n_\theta, n_\phi)$: 3-1-Spin-Index.
\end{itemize}
Beide Sätze von Wicklungszahlen leben auf derselben kompakten
Mannigfaltigkeit. Sie sind nicht unabhängig.

\subsection{Hilbertraum als Spinorbündel}

Mathematisch korrekt ist der FFGFT-Hilbertraum für ein Teilchen
mit Spin auf $T^3$ \emph{kein} Tensorprodukt, sondern der Raum der
Schnitte eines Spinorbündels:
\begin{equation}
  \boxed{
    \mathcal{H}^{T^3}_{\text{FFGFT}}
    = L^2(T^3, S),
  }
\end{equation}
wobei $S$ das Spinorbündel über $T^3$ ist. In dieser Struktur sind
Bahn- und Spin-Operatoren \emph{nicht mehr unabhängig}:
\begin{equation}
  [L_i, S_j] \neq 0
  \qquad \text{mit Korrekturen} \sim \xi.
\end{equation}

\subsection{Konsequenz: geometrische Spin-Bahn-Kopplung}

Die Spin-Bahn-Kopplung ist in FFGFT \emph{nicht} ein zusätzlicher
Hamilton-Term, sondern eine \textbf{geometrische Eigenschaft des
Bündels}. Die \enquote{Spin-Bahn-Konstante} $\xi(r)$ ist nicht
ein freier Parameter, sondern folgt aus der Wicklungsgeometrie.
Insbesondere: das Verhalten von $\xi(r) \propto Z^4 \alpha^2$ in
schweren Atomen folgt aus der spezifischen Wicklungs-Topologie
(Dok.~174).

\subsection{Mathematisches Vorbild}

Die Bündelstruktur ist Standard-Differentialgeometrie -- jedes
moderne Differentialgeometrie-Lehrbuch behandelt Spinorbündel.
Der spezifische FFGFT-Schritt ist die geometrische Identifikation:
\enquote{Spin und Ort sind nicht zwei Hilberträume, sondern zwei
Aspekte derselben Wicklungsstruktur auf einer kompakten
Mannigfaltigkeit.}

\subsection{Was diese Erweiterung kostet}

\begin{itemize}[leftmargin=2em,itemsep=2pt,topsep=2pt]
  \item Hilbertraum ist \emph{kein} Tensorprodukt mehr -- die
        algebraische Struktur ist anders.
  \item Pauli-Prinzip wird geometrisch: zwei Fermionen können nicht
        denselben Wicklungs-Zustand besetzen, weil die
        Bündel-Geometrie das verbietet (Dok.~174 §3).
  \item Drehimpuls-Algebra ist mit Translations-Algebra verschränkt.
  \item Konsequenz: alle phänomenologischen
        \enquote{Spin-Bahn-Konstanten} der Atomphysik sind aus
        FFGFT herleitbar -- sie sind nicht frei.
\end{itemize}

% ==============================================================================
\section{\texorpdfstring{Erweiterung 4: zeitliche Wicklung $k_t$}{Erweiterung 4: zeitliche Wicklung k_t}}
% ==============================================================================

\subsection{Was die Standard-QM mit der Zeit macht}

In der Standard-QM ist die Zeit eine \emph{kontinuierliche, lineare}
Achse $\mathbb{R}$. Die Schrödinger-Gleichung:
\begin{equation}
  i\hbar \frac{\partial \psi}{\partial t} = H \psi,
  \qquad
  E = \hbar \omega
  \qquad \text{(Standard-Dispersion)}.
\end{equation}
Die Energie ist linear in der Frequenz $\omega = 2\pi/T$.

\subsection{FFGFT: Zeit auf $S^1$}

In FFGFT ist die Zeit eine \emph{kompakte} Wicklungsrichtung mit
Quantenzahl $k_t \in \mathbb{Z}$ (Dok.~210). Die Energie wird:
\begin{equation}
  \boxed{
    E = \hbar f(k_t),
  }
\end{equation}
wobei $f$ eine periodische Funktion ist. Im niedrigsten
Wicklungssektor reduziert sich das auf $E \approx \hbar \omega$, aber
die Wicklungsstruktur ist real und trägt physikalische Bedeutung.

\subsection{Mass aus Wicklung}

Die zentrale FFGFT-Aussage: \textbf{die Masse einer Mode ist die
$k_t$-Wicklungsenergie}:
\begin{equation}
  m_i c^2 = \hbar f(k_t^{(i)})
  \qquad \text{für die Mode $i$}.
\end{equation}
In Kombination mit der Yukawa-Form (Dok.~006):
\begin{equation}
  m_i = r_i \cdot \xi_0^{p_i} \cdot v,
\end{equation}
folgen alle 12 Standardmodell-Fermionenmassen aus den Tupeln
$(n_x, n_y, n_z, k_t^{(i)})$ ohne freie Parameter.

\subsection{Hilbertraum-Erweiterung}

Statt $\mathcal{H} = L^2(\mathbb{R}^3) \otimes \mathbb{C}^2$ und Zeit
als kontinuierlicher Parameter, wird der volle FFGFT-Hilbertraum:
\begin{equation}
  \mathcal{H}^{T^4}_{\text{FFGFT}} = L^2(T^4, S),
\end{equation}
mit Schnitten eines Spinorbündels über dem 4-Torus. Die
Schrödinger-Gleichung wird zu einer \emph{Mode-Gleichung} auf $T^4$:
\begin{equation}
  i\hbar \frac{\partial \varphi_{n,l,j}}{\partial t}
  = \mathcal{D}_{T^4} \varphi_{n,l,j},
\end{equation}
wobei $\mathcal{D}_{T^4}$ ein Dirac-artiger Operator auf dem 4-Torus
ist. Im Niederenergielimit reduziert sich dies auf die
Standard-Schrödinger-Gleichung mit fester Masse.

\subsection{Mathematisches Vorbild}

Kompakte Extra-Dimensionen sind seit Kaluza (1921) und Klein (1926)
ein etablierter Rahmen. FFGFT verwendet diese Struktur mit zwei
Modifikationen:
\begin{itemize}[leftmargin=2em,itemsep=2pt,topsep=2pt]
  \item Die Extra-Dimension ist \emph{zeitlich}, nicht räumlich.
  \item Die Wicklungs-Quantenzahl $k_t$ ist nicht frei, sondern
        durch Geometrie und Higgs-Matching festgelegt.
\end{itemize}

\subsection{Was diese Erweiterung kostet}

\begin{itemize}[leftmargin=2em,itemsep=2pt,topsep=2pt]
  \item Schrödinger-Gleichung selbst wird modifiziert.
  \item Massen sind nicht mehr externe Parameter.
  \item Die Standard-QFT-Konstruktion (Lagrangedichten mit
        Yukawa-Termen) ist eine \emph{Niederenergie-Approximation}
        der vollen FFGFT-Mode-Theorie.
  \item Konsequenz: alle 12 Fermionenmassen, die
        Feinstrukturkonstante, und die Higgs-Selbstkopplung
        folgen parameterfrei aus der 4D-Wicklungsgeometrie.
\end{itemize}

% ==============================================================================
\section{Die strukturelle Pointe: bekannte Mathematik, neue Kombination}
% ==============================================================================

Jede der vier Erweiterungen entspricht einer mathematisch \emph{wohlbekannten} Struktur:

\begin{center}
\renewcommand{\arraystretch}{1.4}
{\small
\adjustbox{max width=\textwidth}{%
\begin{tabular}{>{\raggedright\arraybackslash}p{1.0cm}>{\raggedright\arraybackslash}p{3.5cm}>{\raggedright\arraybackslash}p{5.5cm}}
\toprule
\textbf{Erw.} & \textbf{Mathematik} & \textbf{Etabliert seit} \\
\midrule
1 & Quantengruppen $SU_q(2)$ & Drinfeld/Jimbo 1985 \\
2 & Wicklungs-Anregungen & Stringtheorie, Kac-Moody (1980er) \\
3 & Spinorbündel & Differentialgeometrie 1950er \\
4 & Kompakte Extradimension & Kaluza 1921, Klein 1926 \\
\bottomrule
\end{tabular}%
}
}
\renewcommand{\arraystretch}{1.0}
\end{center}

Keine dieser Strukturen ist neu. Was FFGFT \emph{spezifisch} macht,
ist die \textbf{Kombination}: alle vier Erweiterungen gemeinsam und
mit einer einzigen geometrischen Begründung.

\textbf{Was keine andere Theorie tut:}
\begin{itemize}[leftmargin=2em,itemsep=2pt,topsep=2pt]
  \item Quantengruppen werden in der Mathematik studiert, aber nicht
        durch Geometrie aus einer Raumzeit-Dimension begründet.
  \item Stringtheorie hat Wicklungs-Anregungen, aber benötigt 10 oder
        11 Dimensionen und liefert die Standardmodell-Massen nicht
        parameterfrei.
  \item Spinorbündel sind Standardgeometrie, aber die typischen QFT-
        Anwendungen geben Spin und Ort als Tensorprodukt vor.
  \item Kaluza-Klein-Theorien existieren in zahlreichen Varianten,
        aber keine liefert alle 12 SM-Massen aus einer einzigen
        kompakten Dimension.
\end{itemize}

FFGFT ist die \emph{Konvergenz} dieser vier Erweiterungen mit einer
gemeinsamen geometrischen Quelle ($T^4$ und $\xi_0$).

% ==============================================================================
\section{Was eine Hilbertraum-erweiterte Formulierung leistet}
% ==============================================================================

Wenn man den Standard-Hilbertraum mit allen vier Erweiterungen
versieht, erhält man eine vollständige FFGFT-äquivalente
Formulierung, die für Mathematiker und QM-Praktiker zugänglich ist:

\begin{equation}
  \mathcal{H}^{\text{FFGFT}}
  = L^2(T^4, S) \;\bigotimes\; \text{(deformierte $SU_q(2)$-Algebra)},
\end{equation}
mit zusätzlicher Wicklungs-Quantenzahl $n$ pro
Spin-Drehfreiheitsgrad und einer modifizierten Mode-Gleichung
$\mathcal{D}_{T^4}\varphi = E \varphi$ statt der
Schrödinger-Gleichung.

In dieser Formulierung sind alle FFGFT-Vorhersagen direkt aus der
Hilbertraum-Algebra ableitbar:
\begin{itemize}[leftmargin=2em,itemsep=2pt,topsep=2pt]
  \item Massen aus den $k_t$-Wicklungen,
  \item Feinstrukturkonstante aus dem $SU_q(2)$-Casimir,
  \item Spin-Bahn-Kopplung aus der Bündelstruktur,
  \item $K_{\text{frak}}$-Korrekturen aus der Quantengruppen-Deformation,
  \item Casimir-Effekt und kosmologischer Energieverlust aus den
        Wicklungs-Vakuum-Strukturen.
\end{itemize}

\textbf{Methodische Konsequenz:} FFGFT lässt sich \emph{vollständig
in Hilbertraum-Sprache formulieren}, vorausgesetzt man akzeptiert die
vier Erweiterungen. Wer den Standard-Hilbertraum minimal erweitern
möchte, kann mit Stufe 1 ($SU_q(2)$) beginnen, dann Stufe 2
(Wicklungs-Quantenzahl), und so weiter -- jede Stufe bringt eine
weitere Klasse von FFGFT-Aussagen ins Hilbertraum-Vokabular.

% ==============================================================================
\section{Zusammenfassung}
% ==============================================================================

Der Standard-Hilbertraum ($\mathbb{C}^2$ pro Qubit, $L^2(\mathbb{R}^3)$
pro Teilchen, Tensorprodukt für zusammengesetzte Systeme) muss durch
\textbf{vier Erweiterungen} ergänzt werden, um die volle FFGFT
abzubilden:

\begin{enumerate}[leftmargin=2em,itemsep=2pt,topsep=2pt]
  \item \textbf{Deformierte Algebra} $SU_q(2)$ mit
        $q = e^{i\pi(1-K_{\text{frak}})}$.
  \item \textbf{Zyklische Wicklungs-Quantenzahl} $n \in \mathbb{Z}$ für jede
        kompakte Dimension; Hilbertraum wird
        $\bigoplus_n \mathbb{C}^2_n$.
  \item \textbf{Bündelstruktur} statt Tensorprodukt zwischen Ort und
        Spin: $L^2(T^3, S)$ mit gekoppelten Operatoren.
  \item \textbf{Zeitliche Wicklungsdimension} $k_t$ mit Energie
        $E = \hbar f(k_t)$ statt $E = \hbar\omega$; Massen aus
        $k_t$-Wicklungen.
\end{enumerate}

Jede Erweiterung entspricht einer etablierten mathematischen
Struktur (Quantengruppen, Stringtheorie-Wicklungen, Differentialgeometrie,
Kaluza-Klein). FFGFT ist die spezifische, geometrisch begründete
Kombination dieser vier Strukturen, die parameterfrei alle
Standardmodell-Konstanten liefert.

\textbf{Methodische Position:}
\begin{itemize}[leftmargin=2em,itemsep=2pt,topsep=2pt]
  \item FFGFT \emph{fügt sich ein} in den Hilbertraum-Formalismus,
        sofern dieser entsprechend erweitert wird.
  \item Die Erweiterungen sind \emph{nicht ad hoc}, sondern folgen
        aus der spezifischen geometrischen Struktur des $T^4$.
  \item Mathematiker und QM-Praktiker können FFGFT \emph{Stufe für
        Stufe} adaptieren -- mit jeder Erweiterung erhalten sie eine
        weitere Klasse von vorher externen Parametern als
        herleitbare Größen.
\end{itemize}

\textbf{Verbindung zu Dok.~230:} während Dok.~230 die
\enquote{Abwärts-Übersetzung} (FFGFT zum Standard-Hilbertraum, mit
Reduktionen) beschrieb, formuliert dieses Dokument die
\enquote{Aufwärts-Übersetzung} (Hilbertraum-Erweiterungen, um FFGFT
nativ zu tragen). Beide Richtungen sind notwendig, um die
strukturelle Verbindung vollständig zu verstehen.
